% GENERATED FILE. Edit canonical lessons, then run npm run generate:books.
% canonical-source-hash: fnv1a64:d3b50f26a56d0ce1
% canonical-lessons: ML-C204-pavam, ML-C204-alukkulla, ML-C204-minusamulla, ML-C204-parukkan, ML-C204-urunta

\chapter{Poor, Dirty, Smooth, Rough, Round}
\label{ch:ml-poor-dirty-smooth-rough-round}
\hlchaptermodality{204}

\begin{chapteropening}
\textbf{By the end of this chapter:} \emph{I can say poor, dirty, smooth, rough and round.}

Complete the last lesson of chapter 204: I can say poor, dirty, smooth, rough and round.
\end{chapteropening}

\section[\texorpdfstring{\ml{പാവം} (pāvaṁ)}{പാവം (pāvaṁ)}]{\ml{പാവം} (pāvaṁ) — poor, innocent, harmless (of a person)}

\label{lesson:ML-C204-pavam}



\begin{quote}
Before the new one: say the Malayalam for tight, then the Malayalam for empty.
\end{quote}

\subsection*{You'll want to know: \ml{പാവം}}
\textbf{\ml{പാവം}} — \emph{pāvaṁ} — \textquotedblleft{}poor, innocent, harmless (of a person)\textquotedblright{}.

\begin{cousinweb}
\textbf{\ml{പാവം}} (\emph{pāvaṁ}) is usually traced to Sanskrit \emph{pāpa}, \textquotedblleft{}wicked; wretched\textquotedblright{} (CDIAL 8092); in Malayalam the sense softened to \textquotedblleft{}poor, harmless\textquotedblright{}.
\end{cousinweb}

\begin{grammarlens}[title={What you've built}]
One more word for this chapter.
\end{grammarlens}

\subsection*{Your turn}
\begin{itemize}
\raggedright
  \item \emph{Say it:} \emph{pāvaṁ}
  \item \emph{Say it:} \emph{pāvaṁ}, once more
  \item \emph{Say it:} say \emph{pāvaṁ}
  \item \emph{Recall:} say the Malayalam for tight, then the Malayalam for empty, then say \emph{pāvaṁ} again
\end{itemize}

\begin{tcolorbox}[breakable,colback=teal!4,colframe=teal!35!black,title={Before you move on}]
What does \ml{പാവം} mean? (\textquotedblleft{}Poor, innocent, harmless (of a person)\textquotedblright{}.) Say it once more.
\end{tcolorbox}
\section[\texorpdfstring{\ml{അഴുക്കുള്ള} (aḻukkuḷḷa)}{അഴുക്കുള്ള (aḻukkuḷḷa)}]{\ml{അഴുക്കുള്ള} (aḻukkuḷḷa) — dirty}

\label{lesson:ML-C204-alukkulla}



\begin{quote}
Before the new one: say the Malayalam for empty, then the Malayalam for poor.
\end{quote}

\subsection*{You'll want to know: \ml{അഴുക്കുള്ള}}
\textbf{\ml{അഴുക്കുള്ള}} — \emph{aḻukkuḷḷa} — \textquotedblleft{}dirty\textquotedblright{}.

\begin{cousinweb}
\textbf{\ml{അഴുക്കുള്ള}} (\emph{aḻukkuḷḷa}) is native Dravidian, \textquotedblleft{}having dirt\textquotedblright{}: \emph{aḻukku}, \textquotedblleft{}dirt, filth\textquotedblright{} (DEDR 283, with Tamil \emph{aḻukku}), with \emph{-uḷḷa}, \textquotedblleft{}that has\textquotedblright{} (DEDR 697).
\end{cousinweb}

\begin{grammarlens}[title={What you've built}]
One more word for this chapter.
\end{grammarlens}

\subsection*{Your turn}
\begin{itemize}
\raggedright
  \item \emph{Say it:} \emph{aḻukkuḷḷa}
  \item \emph{Say it:} \emph{aḻukkuḷḷa}, once more
  \item \emph{Say it:} say \emph{aḻukkuḷḷa}
  \item \emph{Recall:} say the Malayalam for empty, then the Malayalam for poor, then say \emph{aḻukkuḷḷa} again
\end{itemize}

\begin{tcolorbox}[breakable,colback=teal!4,colframe=teal!35!black,title={Before you move on}]
What does \ml{അഴുക്കുള്ള} mean? (\textquotedblleft{}Dirty\textquotedblright{}.) Say it once more.
\end{tcolorbox}
\section[\texorpdfstring{\ml{മിനുസമുള്ള} (minusamuḷḷa)}{മിനുസമുള്ള (minusamuḷḷa)}]{\ml{മിനുസമുള്ള} (minusamuḷḷa) — smooth}

\label{lesson:ML-C204-minusamulla}



\begin{quote}
Before the new one: say the Malayalam for poor, then the Malayalam for dirty.
\end{quote}

\subsection*{You'll want to know: \ml{മിനുസമുള്ള}}
\textbf{\ml{മിനുസമുള്ള}} — \emph{minusamuḷḷa} — \textquotedblleft{}smooth\textquotedblright{}.

\begin{cousinweb}
\textbf{\ml{മിനുസമുള്ള}} (\emph{minusamuḷḷa}) is \textquotedblleft{}having gloss\textquotedblright{}: \emph{minusam}, \textquotedblleft{}smoothness\textquotedblright{}, most likely from the verb \textquotedblleft{}to glitter, be polished\textquotedblright{}; DEDR 4876 lists Malayalam \emph{minukkam}, \textquotedblleft{}shining, polish\textquotedblright{}, with Tamil \emph{miṉ}, though not \emph{minusam} itself. Then comes \emph{-uḷḷa}, \textquotedblleft{}that has\textquotedblright{} (DEDR 697).
\end{cousinweb}

\begin{grammarlens}[title={What you've built}]
One more word for this chapter.
\end{grammarlens}

\subsection*{Your turn}
\begin{itemize}
\raggedright
  \item \emph{Say it:} \emph{minusamuḷḷa}
  \item \emph{Say it:} \emph{minusamuḷḷa}, once more
  \item \emph{Say it:} say \emph{minusamuḷḷa}
  \item \emph{Recall:} say the Malayalam for poor, then the Malayalam for dirty, then say \emph{minusamuḷḷa} again
\end{itemize}

\begin{tcolorbox}[breakable,colback=teal!4,colframe=teal!35!black,title={Before you move on}]
What does \ml{മിനുസമുള്ള} mean? (\textquotedblleft{}Smooth\textquotedblright{}.) Say it once more.
\end{tcolorbox}
\section[\texorpdfstring{\ml{പരുക്കൻ} (parukkan)}{പരുക്കൻ (parukkan)}]{\ml{പരുക്കൻ} (parukkan) — rough, coarse}

\label{lesson:ML-C204-parukkan}



\begin{quote}
Before the new one: say the Malayalam for dirty, then the Malayalam for smooth.
\end{quote}

\subsection*{You'll want to know: \ml{പരുക്കൻ}}
\textbf{\ml{പരുക്കൻ}} — \emph{parukkan} — \textquotedblleft{}rough, coarse\textquotedblright{}.

\begin{cousinweb}
\textbf{\ml{പരുക്കൻ}} (\emph{parukkan}) is native Dravidian, built on \emph{parukku}, \textquotedblleft{}rough surface\textquotedblright{}; DEDR 3973 lists it, beside Tamil \emph{paru} and Kannada \emph{parige}. DEDR 3973 compares Sanskrit \emph{paruṣa}, with a query.
\end{cousinweb}

\begin{grammarlens}[title={What you've built}]
One more word for this chapter.
\end{grammarlens}

\subsection*{Your turn}
\begin{itemize}
\raggedright
  \item \emph{Say it:} \emph{parukkan}
  \item \emph{Say it:} \emph{parukkan}, once more
  \item \emph{Say it:} say \emph{parukkan}
  \item \emph{Recall:} say the Malayalam for dirty, then the Malayalam for smooth, then say \emph{parukkan} again
\end{itemize}

\begin{tcolorbox}[breakable,colback=teal!4,colframe=teal!35!black,title={Before you move on}]
What does \ml{പരുക്കൻ} mean? (\textquotedblleft{}Rough, coarse\textquotedblright{}.) Say it once more.
\end{tcolorbox}
\section[\texorpdfstring{\ml{ഉരുണ്ട} (uruṇṭa)}{ഉരുണ്ട (uruṇṭa)}]{\ml{ഉരുണ്ട} (uruṇṭa) — round}

\label{lesson:ML-C204-urunta}



\begin{quote}
Before the new one: say the Malayalam for smooth, then the Malayalam for rough.
\end{quote}

\subsection*{You'll want to know: \ml{ഉരുണ്ട}}
\textbf{\ml{ഉരുണ്ട}} — \emph{uruṇṭa} — \textquotedblleft{}round\textquotedblright{}.

\begin{cousinweb}
\textbf{\ml{ഉരുണ്ട}} (\emph{uruṇṭa}) is native Dravidian, a past form of the verb \textquotedblleft{}to roll\textquotedblright{}: \textquotedblleft{}rolled round\textquotedblright{}; DEDR 664 lists the root with Tamil \emph{uruḷ} and Kannada \emph{uruḷ}. DEDR 664 compares Sanskrit \emph{uṇḍeraka}.
\end{cousinweb}

\begin{grammarlens}[title={What you've built}]
One more word for this chapter.
\end{grammarlens}

\subsection*{Your turn}
\begin{itemize}
\raggedright
  \item \emph{Say it:} \emph{uruṇṭa}
  \item \emph{Say it:} \emph{uruṇṭa}, once more
  \item \emph{Say it:} say \emph{uruṇṭa}
  \item \emph{Recall:} say the Malayalam for smooth, then the Malayalam for rough, then say \emph{uruṇṭa} again
\end{itemize}

\begin{tcolorbox}[breakable,colback=teal!4,colframe=teal!35!black,title={Before you move on}]
What does \ml{ഉരുണ്ട} mean? (\textquotedblleft{}Round\textquotedblright{}.) Say it once more.
\end{tcolorbox}
